\documentclass[11pt,twoside,a4paper]{article}

\usepackage[spanish,activeacute]{babel}
\usepackage[utf8x]{inputenc}
\usepackage{a4}
\usepackage[spanish]{layout}
\setlength{\oddsidemargin}{25mm}
\setlength{\evensidemargin}{25mm}
\setlength{\voffset}{6mm}
\setlength{\textwidth}{330pt}
\setlength{\footskip}{6mm}
\setlength{\textheight}{562pt}
\setlength{\marginparwidth}{35pt}
\setlength{\marginparsep}{10pt}
\setlength{\headheight}{15pt}
\usepackage[T1]{fontenc}
\usepackage{eurosym}
\usepackage{url}
\usepackage{graphicx} % Para incluir imágenes
\usepackage{wrapfig}
\usepackage{caption}
\usepackage{subcaption}
\usepackage{fixltx2e} % para los subíndices
\usepackage{color}
\usepackage{lettrine} % para letras capitales de inicios
\usepackage{setspace}
\usepackage{pdflscape} % Para entorno landscape y páginas en apaisado
\usepackage{draftwatermark}
\SetWatermarkText{\textsc{Borrador}}
\SetWatermarkScale{1.2}
\SetWatermarkLightness{.9}
%%%%%%%%%%%
% fancyhdr %
%%%%%%%%%%%%
\usepackage{fancyhdr}
\pagestyle{fancy}
\fancyhead{} % Borramos las cabeceras por defecto
\fancyhf{} % Borramos todas las cabeceras y estilos por defecto
\fancyhead[RO,LE]{\textbf{Notas sobre \LaTeX}}
\fancyfoot{} % Borramos pies de página por defecto
\fancyfoot[LE,RO]{\thepage}
\fancyfoot[CE,CO]{Adolfo Antón Bravo}
\fancyfoot[RE,LO]{Curso de \LaTeX}

\renewcommand{\headrulewidth}{0.4pt}
\renewcommand{\footrulewidth}{0.4pt}
\begin{document}
\title{Formatos en \LaTeX}
\author{Adolfo Antón Bravo}
\maketitle{}
\section{El formato de fuente}
\label{sec:formato}

\subsection{Caracteres especiales}
\label{sec:caracteres}
\LaTeX distingue entre mayúsculas y minúsculas cuando escribimos intrucciones, así que hay que tener cuidado al escribirlas. Por ejemplo, \textbackslash LaTeX sirve para insertar el logotipo de \LaTeX en nuestros documentos.
También podemos escribir ``entre comillas'' -aunque al principio resulta un poco extraño-. Existen ademaás muchos otros símbolos que podemos añadir: \%, \$, \copyright, \ldots

\subsection{Espacios y saltos de línea}
\label{sec:espacioslinea}

También es posible añadir espacios \hspace{2cm} entre palabras, aunque no suele utilizarse demasiado.
Para saltar de párrafo utilizamos la orden \texttt{\textbackslash par} o simplemente dejamos una línea en blanco.

\subsection{El tamaño de la fuente}
\label{sec:fuente}

También podemos utilizar distintos tamaños de fuente que siempre dependerán de la clase del documento y de su configuración.

Existe letra \tiny minúscula, \scriptsize tamaño nota \footnotesize pie de página, \small pequeña, 
ormalsize normal, \large grande, \Large más grande, \LARGE todavía más grande, \huge enorme y \Huge súper enorme. 

\subsection{El estilo de la fuente}
\label{sec:estilo}

ormalsize
Además, podemos escribir en \textbf{negrita}, \textit{cursiva}, \emph{enfatizado}, \uppercase{todo en mayúscula}, \textsc{capitalizado o versales}, \texttt{tamaño fijo}, \ldots

\subsection{Subíndices y superíndices}
\label{sec:subysuper}

La fórmula del ácido sulfúrico es H\textsubscript{2}SO\textsubscript{4}. Y la famosa fórmula de Einstein es \textit{E=mc\textsuperscript{2}}

\subsection{Color}
\label{sec:color}
Para escribir con color debemos utilizar el paquete \color{red}\huge color

ormalsize

\color{black}
\section{El formato de párrafo}
\label{sec:formatoparrafo}

\subsection{Cajas de texto}
\label{sec:cajastexto}
\mbox{Mi texto normal}

\makebox[10cm][r]{Mi texto alineado a la derecha}

\makebox[10cm][s]{Mi texto estirado para ocupar la caja}

\framebox[5cm][c]{Texto centrado con borde}

\framebox[10cm][s]{Una caja que ocupa exactamente lo que ocupa el texto}

\framebox[1.5\width][c]{Una caja que ocupa 1,5 veces más que lo que ocupa el texto}

\framebox[0.5\width][l]{Una caja que ocupa la mitad (0.5) de lo que ocupa el texto}

\framebox[0.3\linewidth][c]{Una caja que ocupa un tercio (0.3) de la página}

\subsection{Bordes y sombreado de cajas}
\label{sec:bordesysombreadocajas}
{\fboxsep 8pt \fboxrule 1pt \fbox{
    \parbox[b]{10cm}{Esta es una caja con borde de 1 punto y margen de 8 puntos}
    } % Cierra fbox
} % cierra el grupo

\subsection{Color de las cajas}
\label{sec:colorcajas}

Podemos colocar las cajas con \textbf{fcolorbox}:

\colorbox{red}{texto escrito en negro dentro de una caja de color rojo}
\\

\fcolorbox{blue}{white}{
    \parbox[b]{.2\linewidth}{texto escrito en negro dentro de una caja de párrafo de color negro con borde azul}
}


\subsection{Alineación de las cajas}
\label{sec:alineacioncajas}

Para mostrar la alineación de las cajas utilizaremos cajas de párrafo (\texttt{parbox} con ancho de \texttt{0.3 linewidth}:

\parbox[b]{0.3\linewidth}{\raggedright{Caja de párrafo a la izqda. con texto alineado a la izqda.}}
\parbox[b]{0.3\linewidth}{\centering Caja de párrafo centrada con texto centrado}
\parbox[b]{0.3\linewidth}{\raggedleft{Caja de párrafo a la dcha. con texto alineado a la dcha.}}

\subsection{Espaciado entre párrafos}
\label{sec:espaciadoparrafos}

Podemos saltar de párrafo dejando una línea en blanco con \texttt{par}\par Y de página con con \texttt{newpage}
ewpage{}
Ésta es otra página. y también podemos saltar de línea sin saltar de párrafo con \texttt{newline} o dos barras invertidas.
ewline{} Esto es un salto de línea. Puede observarse que en los saltos de línea no se indenta mientras que en los saltos de párrafo, sí.

Entre los párrafos puede haber espacios pequeños:\smallskip{}

Espacios medios:
\medskip{}

O espacios grandes:
\bigskip{}

También podemos dejar el espacio que queramos con \texttt{vspace}\vspace{4cm}

Acabamos de dejar 4 cm. de espacio.
\subsection{Letras capitales}
\label{sec:letrascapitales}
\lettrine{E}{n un lugar de la Mancha}, de cuyo nombre no quiero acordarme, no ha mucho que vivía un hidalgo de los de lanza en astillero\ldots{}

\subsection{Interlineado}
\label{sec:interlineado}

Normalmente hay que dejar a \LaTeX{} realizar el interlineado dependiendo del tipo de documento.

\doublespacing
Pero también podemos modificarlo con el paquete \texttt{setspace}. Por ejemplo, este párrafo tiene espaciado doble.

\singlespacing
Ahora volvemos al espaciado normal mediante la instrucción \texttt{singlespacing}. A partir de ahora el resto del texto tiene espaciado normal.

\section{Entornos y listas}
\label{sec:entornosylistas}

Modificar el formato de párrafo directamente es muy complicado. Mejor utilizar entornos.

\subsection{Entornos de alineación}
\label{sec:entornoalineacion}

El texto que escribimos dentro de este entorno está alineado a la izquierda y queda irregulas por la derecha. Como se puede ver, es distinto a la alineación justificada que \LaTeX{} utiliza por defecto.

\begin{flushright}
  El texto que escribimos dentro de este entorno está alineado por la derecha, por lo que queda irregular por la izquierda.
\end{flushright}

\begin{center}
  Este texto queda centrado.\\
Como siempre, podemos cambiar de líena sin cambiar de párrafo utilizando las dos barras invertidas.
\end{center}

\subsection{Entorno para crear cajas de párrafos}
\label{sec:entornocajasparrafos}
Mucho más sencillo que \texttt{parbox}:

\begin{minipage}[l]{0.4\linewidth}
Esta es una caja de texto que ocupa el 40\%
\end{minipage}
\hfill{}
\begin{minipage}[r]{0.4\linewidth}
  Esta es otra caja de texto que ocupa el otro 40\%. Para separarlas podemos utilizar \texttt{hfill}
\end{minipage}

\subsection{Entornos para citas y versos}
\label{sec:entornoscitasyversos}

\subsubsection{Quote}
\label{sec:quote}

Una de las reglas de ortografía de la letra ``b'' dice:
\begin{quote}
  Se escribe ``b'' delante de cualquier consonante
\end{quote}

\subsubsection{Quotation}
\label{sec:quotation}

A la pregunta \textbf{¿cómo ayuda la tecnología móvil a las organizacioens en su objetivo de lograr eficiencia operacional?}, Jorge contesto:
\begin{quotation}
  Les ayuda a hacer primordialmente dos cosas.
En primer lugar, los empleados pueden obtener información de producto e indicadores de rendimiento en tiempo real, en cualquier momento y desde cualquier lugar. Esto acelera el proceso de toma de decisiones.
Y en segundo lugar, la movilidad mejora la comunicación con los empleados, clientes, socios y participantes de la cadena de suministros.
\end{quotation}

\subsubsection{Verse}
\label{sec:verse}

\large Rima XXIII

ormalsize
\begin{verse}
  Por una mirada, un mundo;\\
  por una sonrisa, un cielo;\\
  por un beso\ldots{} yo no sé\\
  qué te diera por un beso.
\end{verse}

\subsection{Verbatim}
\label{sec:verbatim}
\scriptsize
\begin{verbatim}
var sumar = function (a, b) {
 return a + b; // Suma los números
}
var c = sumar(2,3);
\end{verbatim}

ormalsize

\subsection{Listas y descripciones}
\label{sec:listasdescripciones}

\subsubsection{Listas con viñetas}
\label{sec:listasviñetas}

Los cuatro puntos cardinales principales son:
\begin{itemize}
\item Norte
\item Sur
\item Este
\item Oeste
\end{itemize}

\subsubsection{Listas numeradas}
\label{sec:listasnumeradas}

Los planetas del Sistema Solar, ordenados por proximidad al Sol, son:
\begin{enumerate}
\item Mercurio
\item Venus
\item Tierra
\item Marte\ldots{}
\end{enumerate}

\subsubsection{Listas de definiciones}
\label{sec:listasdefiniciones}

Una célula animal se compone de las siguientes partes:
\begin{description}
\item[Membrana] Es el límite externo de la célula
\item[Citoplasma] Es el volumen de la célula entre la membrana y el núcleo
\item[Núcleo] Es el orgánulo membranoso que se encuentra en el centro de la célula
\end{description}

ewpage{}
\section{Estructura de página}
\label{sec:estructurapagina}
\layout{}
%\includegraphics{estructura-pagina-izqda.png}

\thispagestyle{empty}
%\includegraphics{estructura-pagina-dcha.png}
\begin{landscape}
\thispagestyle{empty}
  \subsection{Página en apaisado}
  \label{sec:apaisado}
Esta sección la veremos apaisada
  
\end{landscape}

ewpage
Aquí volvemos a la normalidad

\end{document}

